\chapter{Resultados e discussão}
Apresente os resultados obtidos e discuta seu significado à luz da literatura revisada. A Tabela~\ref{tab:exemplo-csv} demonstra como importar dados diretamente de um arquivo CSV externo utilizando o pacote \texttt{csvsimple}.

% ATENÇÃO: o arquivo CSV deve estar previamente tratado e pronto para exibição
% (casas decimais arredondadas, separador decimal com vírgula no padrão
% brasileiro, separador de colunas ";" e eventuais totais já calculados).
% O LaTeX apenas lê e formata a tabela visualmente, sem processar dados brutos.
\begin{table}[htbp]
  \centering
  \caption{Desempenho operacional importado de arquivo CSV}\label{tab:exemplo-csv}
  \tagpdfsetup{table/header-rows={1}}
  \csvreader[
    separator=semicolon,
    respect all,
    tabular=lrrr,
    table head=\toprule \textbf{Linha} & \textbf{Tempo de ciclo (min)} & \textbf{Eficiência (\%)} & \textbf{Produção diária (unid.)} \\ \midrule,
    late after line=\\,
    table foot=\bottomrule
  ]{dados/exemplo-tabela.csv}{}{%
    \csvcoli & \csvcolii & \csvcoliii & \csvcoliv
  }
  \fonte{Elaboração própria}
\end{table}
